\section{End-Term Examination --- Set C}
\label{sec:endterm_set_c}

\LPUPaperHeader{End-Term Examination}{C}{3 Hours}{60 Marks}{CO1 to CO6 (Units 1 to 6 Comprehensive)}

\begin{center}
\sffamily\textbf{\color{AlertRed}Structure of Question Paper:}\\[1mm]
\textbf{Part A (Objective):} 30 MCQs strictly from \textbf{Units 4, 5, and 6}. Total weight: \textbf{20 Marks}. All questions compulsory.\\
\textbf{Part B (Subjective):} 6 Questions from \textbf{all units} (10 Marks each). \textbf{Attempt any 4 questions} ($4 \times 10 = \mathbf{40\text{ Marks}}$).\\
\textbf{Total Maximum Marks: 60 Marks.}
\end{center}

\vspace{3mm}

% ------------------------------------------------------------------------------
% PART A: 30 MCQS (UNITS 4, 5, 6)
% ------------------------------------------------------------------------------
\subsection*{PART A: Objective Section (30 MCQs --- Units 4, 5, 6) [20 Marks]}

\mcqitem{1}{The value of $\lim_{(x,y)\to(0,0)} \frac{x^4 y^4}{(x^2 + y^4)^3}$ along the path $x = y^2$ is:}{$1/8$}{$1/4$}{$0$}{$1$}{1}

\mcqitem{2}{If $u = \tan^{-1}\left(\frac{x^3 + y^3}{x - y}\right)$, then $x \frac{\partial u}{\partial x} + y \frac{\partial u}{\partial y}$ is equal to:}{$\sin 2u$}{$\cos 2u$}{$2\tan u$}{$\tan 2u$}{1}

\mcqitem{3}{The saddle point of $f(x,y) = 2x^2 - xy - 3y^2 - 3x + 7y$ is:}{$(1, 1)$}{$(0, 0)$}{$(1, -1)$}{$(2, 1)$}{1}

\mcqitem{4}{The directional derivative of $\phi = 2xy + z^2$ at $(1, -1, 3)$ in the direction of $2\hat{i} + 2\hat{j} - \hat{k}$ is:}{$-14/3$}{$-10/3$}{$14/3$}{$4$}{1}

\mcqitem{5}{If $\vec{F} = (x+2y)\hat{i} + (y+3z)\hat{j} + (z+cx)\hat{k}$ is solenoidal, then $c$ can be:}{Any value because $\nabla \cdot \vec{F}$ is independent of $c$}{$0$}{$-1$}{$-3$}{1}

\mcqitem{6}{The curl of the vector field $\vec{F} = x^2\hat{i} + y^2\hat{j} + z^2\hat{k}$ is:}{$\vec{0}$}{$2x\hat{i} + 2y\hat{j} + 2z\hat{k}$}{$2(x+y+z)$}{$\vec{r}$}{1}

\mcqitem{7}{The value of $\int_0^1 \int_0^x (x^2 + y^2) \, dy \, dx$ is:}{$1/3$}{$1/6$}{$1/4$}{$5/12$}{1}

\mcqitem{8}{In changing from Cartesian to spherical coordinates, the Jacobian determinant $|J|$ is:}{$\rho^2 \sin\phi$}{$\rho \sin\phi$}{$\rho^2 \cos\phi$}{$\rho \cos\phi$}{1}

\mcqitem{9}{If $\nabla \times \vec{F} = \vec{0}$, the line integral $\int_C \vec{F} \cdot d\vec{r}$ is:}{Path independent}{Path dependent}{$0$ always}{$\infty$}{1}

\mcqitem{10}{The value of $\nabla^2(r^3)$ where $r = \|\vec{r}\|$ is:}{$12r$}{$6r$}{$3r$}{$12r^2$}{1}

\mcqitem{11}{The value of $\oint_C \vec{r} \cdot d\vec{r}$ around any closed smooth curve $C$ is:}{$0$}{$2\pi$}{$\pi$}{Enclosed area}{1}

\mcqitem{12}{The function $f(x,y) = x^2 - y^2$ has at $(0,0)$:}{A local minimum}{A local maximum}{A saddle point}{No critical point}{1}

\mcqitem{13}{The value of $\int_0^2 \int_0^1 (xy + e^x) \, dy \, dx$ is:}{$e^2 + 1$}{$e^2 - 1$}{$e^2$}{$2e^2$}{1}

\mcqitem{14}{A surface integral $\iint_S \vec{F} \cdot \hat{n} \, dS$ represents the:}{Total flux of $\vec{F}$ through $S$}{Work done}{Circulation}{Volume}{1}

\mcqitem{15}{If $\vec{F} = y\hat{i} - x\hat{j}$, then $\nabla \times \vec{F}$ is:}{$-2\hat{k}$}{$2\hat{k}$}{$\vec{0}$}{$\hat{k}$}{1}

\mcqitem{16}{The area enclosed by the astroid $x^{2/3} + y^{2/3} = a^{2/3}$ is given by:}{$\frac{3}{8}\pi a^2$}{$\frac{3}{4}\pi a^2$}{$\pi a^2$}{$\frac{1}{2}\pi a^2$}{1}

\mcqitem{17}{The value of $\iint_S \hat{n} \, dS$ over any closed surface $S$ is identically:}{$\vec{0}$}{$1$}{Surface Area}{$V\hat{k}$}{1}

\mcqitem{18}{The stationary value of $x^2 + y^2 + z^2$ subject to $x + y + z = 3$ is:}{$3$}{$9$}{$1$}{$6$}{1}

\mcqitem{19}{If $u, v$ are functionally dependent, their Jacobian $\frac{\partial(u,v)}{\partial(x,y)}$ is:}{$0$}{$1$}{$-1$}{$\infty$}{1}

\mcqitem{20}{The volume of the ellipsoid $\frac{x^2}{a^2} + \frac{y^2}{b^2} + \frac{z^2}{c^2} \le 1$ is:}{$\frac{4}{3}\pi abc$}{$4\pi abc$}{$\frac{2}{3}\pi abc$}{$\pi abc$}{1}

\mcqitem{21}{The value of $\nabla(\vec{a} \cdot \vec{r})$ where $\vec{a}$ is a constant vector is:}{$\vec{a}$}{$\vec{0}$}{$2\vec{a}$}{$\vec{r}$}{1}

\mcqitem{22}{The value of the integral $\int_0^1 \int_0^{\sqrt{1-x^2}} dy \, dx$ is:}{$\frac{\pi}{4}$}{$\frac{\pi}{2}$}{$\pi$}{$1$}{1}

\mcqitem{23}{If $f(x,y)$ has continuous partial derivatives, then $\frac{\partial^2 f}{\partial x \partial y} = \frac{\partial^2 f}{\partial y \partial x}$ is guaranteed by:}{Schwarz's Theorem}{Euler's Theorem}{Taylor's Theorem}{Green's Theorem}{1}

\mcqitem{24}{The divergence of $\vec{r}/r^3$ for $r \neq 0$ is:}{$0$}{$3/r^3$}{$-3/r^4$}{$1/r^2$}{1}

\mcqitem{25}{The double integral $\iint_R dx \, dy$ over $0 \le x \le 2, 0 \le y \le 3$ is:}{$6$}{$5$}{$1$}{$12$}{1}

\mcqitem{26}{The unit outward normal to the sphere $x^2 + y^2 + z^2 = R^2$ at any point is:}{$\frac{\vec{r}}{R}$}{$\frac{\vec{r}}{R^2}$}{$\frac{\nabla r}{R}$}{$\vec{r}$}{1}

\mcqitem{27}{The value of $\oint_C \nabla \phi \cdot d\vec{r}$ around any closed curve is:}{$0$}{$\phi$}{$2\pi$}{$1$}{1}

\mcqitem{28}{If $u = x+y+z, v = x y + y z + z x, w = x^2 + y^2 + z^2$, the Jacobian $\frac{\partial(u,v,w)}{\partial(x,y,z)}$ is:}{$0$}{$1$}{$2$}{$6$}{1}

\mcqitem{29}{The volume of a right circular cone of base radius $R$ and height $h$ is:}{$\frac{1}{3}\pi R^2 h$}{$\pi R^2 h$}{$\frac{2}{3}\pi R^2 h$}{$\frac{1}{2}\pi R^2 h$}{1}

\mcqitem{30}{If $\vec{F} = \text{curl}(\vec{A})$, then the flux of $\vec{F}$ across any closed surface $S$ is:}{$0$}{$1$}{Area of $S$}{Volume of $S$}{1}

\vspace{5mm}
\hrule
\vspace{5mm}

% ------------------------------------------------------------------------------
% PART B: 6 SUBJECTIVE QUESTIONS (ATTEMPT ANY 4)
% ------------------------------------------------------------------------------
\subsection*{PART B: Descriptive Section (Attempt Any 4 of 6) [40 Marks]}

\questionitem{31 (Unit 1)}{
    Diagonalize the real symmetric matrix $A$ by an orthogonal transformation $P$:
    \[
    A = \begin{bmatrix} 3 & -1 & 1 \\ -1 & 5 & -1 \\ 1 & -1 & 3 \end{bmatrix}
    \]
    Find the modal matrix $P$ and the resulting diagonal matrix $D = P^T A P$.
}{10}{CO1}{L3 Apply}

\questionitem{32 (Unit 2)}{
    Using the \textbf{Method of Variation of Parameters}, solve the second-order ordinary differential equation:
    \[
    \frac{d^2 y}{dx^2} - 2\frac{dy}{dx} + y = \frac{e^x}{x} \quad (x > 0)
    \]
}{10}{CO2}{L4 Analyze}

\questionitem{33 (Unit 3)}{
    \textbf{(a)} Test the convergence of the infinite series:
    \[
    \sum_{n=1}^\infty \frac{n^3 + a}{2^n + b} \quad (a, b > 0)
    \]
    \textbf{(b)} Expand $f(x) = e^x \sin x$ in a \textbf{Maclaurin series} up to the term containing $x^5$.
}{10}{CO3}{L3 Apply}

\questionitem{34 (Unit 4)}{
    Find the maximum and minimum values of the function:
    \[
    f(x,y,z) = x y z
    \]
    subject to the spherical constraint $x^2 + y^2 + z^2 = 1$, using \textbf{Lagrange Multipliers}.
}{10}{CO4}{L3 Apply}

\questionitem{35 (Unit 5)}{
    Evaluate the triple integral:
    \[
    \iiint_V x y z \, dx \, dy \, dz
    \]
    over the \textbf{positive octant} of the ellipsoid:
    \[
    \frac{x^2}{a^2} + \frac{y^2}{b^2} + \frac{z^2}{c^2} \le 1
    \]
}{10}{CO5}{L3 Apply}

\questionitem{36 (Unit 6)}{
    Verify \textbf{Green's Theorem in the plane} for:
    \[
    \oint_C \left[ (xy + y^2) \, dx + x^2 \, dy \right]
    \]
    where $C$ is the closed curve bounding the region between the line $y = x$ and the parabola $y = x^2$.
}{10}{CO6}{L4 Analyze}

\vspace{5mm}
\hrule
\vspace{5mm}

% ------------------------------------------------------------------------------
% COMPLETE STEP-BY-STEP SOLUTIONS
% ------------------------------------------------------------------------------
\subsection*{Solutions \& Marking Scheme (End-Term Set C)}

\subsubsection*{Part A: Answer Key \& Explanations}

\begin{center}
\begin{tabular}{|c|c||c|c||c|c||c|c||c|c||c|c|}
\hline
\textbf{Q} & \textbf{Ans} & \textbf{Q} & \textbf{Ans} & \textbf{Q} & \textbf{Ans} & \textbf{Q} & \textbf{Ans} & \textbf{Q} & \textbf{Ans} & \textbf{Q} & \textbf{Ans} \\
\hline
1 & \textbf{A} & 6 & \textbf{A} & 11 & \textbf{A} & 16 & \textbf{A} & 21 & \textbf{A} & 26 & \textbf{A} \\
2 & \textbf{A} & 7 & \textbf{A} & 12 & \textbf{C} & 17 & \textbf{A} & 22 & \textbf{A} & 27 & \textbf{A} \\
3 & \textbf{A} & 8 & \textbf{A} & 13 & \textbf{A} & 18 & \textbf{A} & 23 & \textbf{A} & 28 & \textbf{A} \\
4 & \textbf{A} & 9 & \textbf{A} & 14 & \textbf{A} & 19 & \textbf{A} & 24 & \textbf{A} & 29 & \textbf{A} \\
5 & \textbf{A} & 10 & \textbf{A} & 15 & \textbf{A} & 20 & \textbf{A} & 25 & \textbf{A} & 30 & \textbf{A} \\
\hline
\end{tabular}
\end{center}

\begin{solutionbox}[Explanations for Part A]
\textbf{Q.1 (Ans: A)} Along $x = y^2$: $\frac{(y^2)^4 y^4}{((y^2)^2 + y^4)^3} = \frac{y^{12}}{(2y^4)^3} = \frac{y^{12}}{8y^{12}} = \frac{1}{8}$.\\[1.5mm]
\textbf{Q.2 (Ans: A)} Let $f = \tan u = \frac{x^3+y^3}{x-y}$ which is homogeneous of degree 2. By Euler's Theorem, $x f_x + y f_y = 2f \implies \sec^2 u (x u_x + y u_y) = 2\tan u \implies x u_x + y u_y = \frac{2\tan u}{\sec^2 u} = 2\sin u\cos u = \sin 2u$.\\[1.5mm]
\textbf{Q.10 (Ans: A)} $\nabla^2(r^n) = n(n+1)r^{n-2}$. For $n=3$: $3(4)r^{3-2} = 12r$.\\[1.5mm]
\textbf{Q.24 (Ans: A)} $\nabla \cdot (\vec{r}/r^3) = \frac{\nabla \cdot \vec{r}}{r^3} + \vec{r} \cdot \nabla(r^{-3}) = \frac{3}{r^3} + \vec{r} \cdot (-3r^{-5}\vec{r}) = \frac{3}{r^3} - \frac{3r^2}{r^5} = 0$.\\[1.5mm]
\textbf{Q.28 (Ans: A)} Since $w = x^2+y^2+z^2 = (x+y+z)^2 - 2(xy+yz+zx) = u^2 - 2v$, $w$ is a functional combination of $u$ and $v$. Hence the Jacobian vanishes identically ($J = 0$).\\[1.5mm]
\textbf{Q.30 (Ans: A)} By Gauss Divergence Theorem, $\iint_S (\nabla \times \vec{A}) \cdot \hat{n} dS = \iiint_V \nabla \cdot (\nabla \times \vec{A}) dV = \iiint_V 0 dV = 0$.
\end{solutionbox}

\subsubsection*{Part B: Detailed Subjective Solutions}

% Solution 31
\begin{solutionbox}[Solution to Question 31 (10 Marks)]
\textbf{Step 1: Eigenvalues of $A$}
$A = \begin{bmatrix} 3 & -1 & 1 \\ -1 & 5 & -1 \\ 1 & -1 & 3 \end{bmatrix}$.
\begin{itemize}[leftmargin=6mm]
    \item $S_1 = 3 + 5 + 3 = 11$.
    \item $S_2 = (15 - 1) + (9 - 1) + (15 - 1) = 14 + 8 + 14 = 36$.
    \item $S_3 = \det(A) = 3(14) - (-1)(-3 + 1) + 1(1 - 5) = 42 - 2 - 4 = 36$.
\end{itemize}
Equation: $\lambda^3 - 11\lambda^2 + 36\lambda - 36 = 0$.
Factorization: $(\lambda - 2)(\lambda - 3)(\lambda - 6) = 0 \implies \mathbf{\lambda_1 = 2, \quad \lambda_2 = 3, \quad \lambda_3 = 6}$.

\textbf{Step 2: Normalized Eigenvectors}
\begin{itemize}[leftmargin=6mm]
    \item For $\lambda_1 = 2$: $(A - 2I)X = 0 \implies \begin{bmatrix} 1 & -1 & 1 \\ -1 & 3 & -1 \\ 1 & -1 & 1 \end{bmatrix}X = 0$.
    $x_1 - x_2 + x_3 = 0$ and $-x_1 + 3x_2 - x_3 = 0 \implies 2x_2 = 0 \implies x_2 = 0, x_3 = -x_1$.
    Eigenvector $X_1 = \begin{bmatrix} 1 \\ 0 \\ -1 \end{bmatrix}$. Normalized: $\mathbf{\hat{X}_1 = \frac{1}{\sqrt{2}}\begin{bmatrix} 1 \\ 0 \\ -1 \end{bmatrix}}$.
    \item For $\lambda_2 = 3$: $(A - 3I)X = 0 \implies \begin{bmatrix} 0 & -1 & 1 \\ -1 & 2 & -1 \\ 1 & -1 & 0 \end{bmatrix}X = 0 \implies x_1 = x_2 = x_3$.
    Eigenvector $X_2 = \begin{bmatrix} 1 \\ 1 \\ 1 \end{bmatrix}$. Normalized: $\mathbf{\hat{X}_2 = \frac{1}{\sqrt{3}}\begin{bmatrix} 1 \\ 1 \\ 1 \end{bmatrix}}$.
    \item For $\lambda_3 = 6$: $(A - 6I)X = 0 \implies \begin{bmatrix} -3 & -1 & 1 \\ -1 & -1 & -1 \\ 1 & -1 & -3 \end{bmatrix}X = 0$.
    Cross-multiplication gives $x_1 = 1, x_2 = -2, x_3 = 1$.
    Eigenvector $X_3 = \begin{bmatrix} 1 \\ -2 \\ 1 \end{bmatrix}$. Normalized: $\mathbf{\hat{X}_3 = \frac{1}{\sqrt{6}}\begin{bmatrix} 1 \\ -2 \\ 1 \end{bmatrix}}$.
\end{itemize}

\textbf{Step 3: Orthogonal Modal Matrix $P$ and Diagonal Matrix $D$}
\[
P = \begin{bmatrix}
1/\sqrt{2} & 1/\sqrt{3} & 1/\sqrt{6} \\
0 & 1/\sqrt{3} & -2/\sqrt{6} \\
-1/\sqrt{2} & 1/\sqrt{3} & 1/\sqrt{6}
\end{bmatrix}
\]
Since $P$ is orthogonal ($P^{-1} = P^T$):
\[
\mathbf{D = P^T A P = \begin{bmatrix} 2 & 0 & 0 \\ 0 & 3 & 0 \\ 0 & 0 & 6 \end{bmatrix}}
\]

\begin{markingrubric}
\textbf{Mark Distribution:}
$\bullet$ Characteristic equation and eigenvalues $\lambda = 2, 3, 6$: \textbf{3 Marks}\\
$\bullet$ Eigenvectors determination: \textbf{3 Marks}\\
$\bullet$ Normalization of eigenvectors: \textbf{2 Marks}\\
$\bullet$ Construction of modal matrix $P$ and diagonal matrix $D$: \textbf{2 Marks}
\end{markingrubric}
\end{solutionbox}

% Solution 32
\begin{solutionbox}[Solution to Question 32 (10 Marks)]
\textbf{Step 1: Complementary Function}
$m^2 - 2m + 1 = 0 \implies (m - 1)^2 = 0 \implies m = 1, 1$.
\[
y_1 = e^x, \quad y_2 = x e^x \implies \mathbf{y_c = (c_1 + c_2 x)e^x}
\]
\textbf{Step 2: Wronskian $W(y_1, y_2)$}
\[
W = \begin{vmatrix} e^x & x e^x \\ e^x & (x+1)e^x \end{vmatrix} = e^{2x}(x+1) - x e^{2x} = e^{2x} \neq 0
\]
\textbf{Step 3: Particular Integral via Variation of Parameters}
Let $y_p = u(x)y_1 + v(x)y_2$ where $R(x) = \frac{e^x}{x}$:
\[
u(x) = -\int \frac{y_2 R(x)}{W} dx = -\int \frac{(x e^x)(e^x / x)}{e^{2x}} dx = -\int 1 \, dx = \mathbf{-x}
\]
\[
v(x) = \int \frac{y_1 R(x)}{W} dx = \int \frac{(e^x)(e^x / x)}{e^{2x}} dx = \int \frac{1}{x} \, dx = \mathbf{\ln x}
\]
\textbf{Step 4: Particular Integral and General Solution}
\[
y_p = (-x)e^x + (\ln x)(x e^x) = x e^x (\ln x - 1)
\]
Absorbing the $-x e^x$ term into the arbitrary constant $c_2 x e^x$ of $y_c$:
\[
\mathbf{y(x) = c_1 e^x + c_2 x e^x + x e^x \ln x}
\]

\begin{markingrubric}
\textbf{Mark Distribution:}
$\bullet$ Complementary function $y_c = (c_1 + c_2 x)e^x$: \textbf{2 Marks}\\
$\bullet$ Wronskian computation $W = e^{2x}$: \textbf{2 Marks}\\
$\bullet$ Integrals $u(x) = -x$ and $v(x) = \ln x$: \textbf{4 Marks}\\
$\bullet$ Final general solution: \textbf{2 Marks}
\end{markingrubric}
\end{solutionbox}

% Solution 33
\begin{solutionbox}[Solution to Question 33 (10 Marks)]
\textbf{Part (a): Convergence of Series (5 Marks)}
Let $u_n = \frac{n^3 + a}{2^n + b}$. As $n \to \infty$, $u_n \approx \frac{n^3}{2^n}$.
Apply D'Alembert's Ratio Test:
\[
L = \lim_{n\to\infty} \frac{u_{n+1}}{u_n} = \lim_{n\to\infty} \frac{(n+1)^3 + a}{2^{n+1} + b} \cdot \frac{2^n + b}{n^3 + a} = \lim_{n\to\infty} \frac{(n+1)^3}{n^3} \cdot \frac{2^n}{2^{n+1}} = (1) \cdot \frac{1}{2} = \mathbf{\frac{1}{2}}
\]
Since $L = \frac{1}{2} < 1$, the series is \textbf{strictly CONVERGENT}.

\textbf{Part (b): Maclaurin Series of $e^x \sin x$ (5 Marks)}
We know the standard Maclaurin series:
\[
e^x = 1 + x + \frac{x^2}{2} + \frac{x^3}{6} + \frac{x^4}{24} + \dots
\]
\[
\sin x = x - \frac{x^3}{6} + \frac{x^5}{120} - \dots
\]
Multiply the series term by term up to degree 5:
\[
e^x \sin x = \left(1 + x + \frac{x^2}{2} + \frac{x^3}{6} + \frac{x^4}{24}\right) \left(x - \frac{x^3}{6} + \frac{x^5}{120}\right)
\]
\begin{itemize}[leftmargin=6mm]
    \item Degree 1: $1 \cdot x = x$
    \item Degree 2: $x \cdot x = x^2$
    \item Degree 3: $1 \cdot (-x^3/6) + \frac{x^2}{2} \cdot x = -\frac{1}{6}x^3 + \frac{1}{2}x^3 = \frac{1}{3}x^3$
    \item Degree 4: $x \cdot (-x^3/6) + \frac{x^3}{6} \cdot x = -\frac{1}{6}x^4 + \frac{1}{6}x^4 = 0$
    \item Degree 5: $1 \cdot (x^5/120) + \frac{x^2}{2}(-x^3/6) + \frac{x^4}{24}(x) = \left(\frac{1}{120} - \frac{1}{12} + \frac{1}{24}\right)x^5 = \frac{1 - 10 + 5}{120}x^5 = -\frac{4}{120}x^5 = -\frac{1}{30}x^5$
\end{itemize}
\[
\mathbf{e^x \sin x = x + x^2 + \frac{1}{3}x^3 - \frac{1}{30}x^5 + \dots}
\]

\begin{markingrubric}
\textbf{Mark Distribution:}
$\bullet$ Part (a) Ratio test limit $L = 1/2 < 1$ (Convergent): \textbf{5 Marks}\\
$\bullet$ Part (b) Series expansion multiplication up to degree 5: \textbf{5 Marks}
\end{markingrubric}
\end{solutionbox}

% Solution 34
\begin{solutionbox}[Solution to Question 34 (10 Marks)]
\textbf{Step 1: Lagrange Equations}
Objective: $f(x,y,z) = xyz$. Constraint: $g(x,y,z) = x^2 + y^2 + z^2 - 1 = 0$.
Set $\nabla f = \lambda \nabla g$:
\begin{align}
yz &= 2\lambda x \label{eq:sph1} \\
xz &= 2\lambda y \label{eq:sph2} \\
xy &= 2\lambda z \label{eq:sph3}
\end{align}
\textbf{Step 2: Solving the System}
Multiply \eqref{eq:sph1} by $x$, \eqref{eq:sph2} by $y$, and \eqref{eq:sph3} by $z$:
\[
xyz = 2\lambda x^2 = 2\lambda y^2 = 2\lambda z^2
\]
If $xyz \neq 0$, then $\lambda \neq 0$, which immediately implies:
\[
x^2 = y^2 = z^2
\]
Substitute into the constraint $x^2 + y^2 + z^2 = 1$:
\[
3x^2 = 1 \implies x^2 = \frac{1}{3} \implies x = \pm \frac{1}{\sqrt{3}}, \quad y = \pm \frac{1}{\sqrt{3}}, \quad z = \pm \frac{1}{\sqrt{3}}
\]
\textbf{Step 3: Extreme Values}
\begin{itemize}[leftmargin=6mm]
    \item When all three coordinates have positive signs (or two negative, one positive):
    \[
    f_{\max} = \left(\frac{1}{\sqrt{3}}\right)\left(\frac{1}{\sqrt{3}}\right)\left(\frac{1}{\sqrt{3}}\right) = \mathbf{\frac{1}{3\sqrt{3}}}
    \]
    \item When all three are negative (or one negative, two positive):
    \[
    f_{\min} = \mathbf{-\frac{1}{3\sqrt{3}}}
    \]
\end{itemize}

\begin{markingrubric}
\textbf{Mark Distribution:}
$\bullet$ Setting up Lagrange system: \textbf{3 Marks}\\
$\bullet$ Deducing $x^2 = y^2 = z^2 = 1/3$: \textbf{4 Marks}\\
$\bullet$ Finding $f_{\max} = \frac{1}{3\sqrt{3}}$ and $f_{\min} = -\frac{1}{3\sqrt{3}}$: \textbf{3 Marks}
\end{markingrubric}
\end{solutionbox}

% Solution 35
\begin{solutionbox}[Solution to Question 35 (10 Marks)]
\textbf{Step 1: Transformation to Unit Sphere}
Substitute:
\[
x = a u, \quad y = b v, \quad z = c w
\]
Then the ellipsoid $\frac{x^2}{a^2} + \frac{y^2}{b^2} + \frac{z^2}{c^2} \le 1$ transforms into the unit sphere:
\[
u^2 + v^2 + w^2 \le 1, \quad u \ge 0, v \ge 0, w \ge 0
\]
The Jacobian of the transformation is:
\[
J = \frac{\partial(x,y,z)}{\partial(u,v,w)} = abc \implies dx \, dy \, dz = abc \, du \, dv \, dw
\]
\textbf{Step 2: Integral in $(u, v, w)$ Space}
\[
I = \iiint_V (a u)(b v)(c w) (abc \, du \, dv \, dw) = a^2 b^2 c^2 \iiint_{V'} u v w \, du \, dv \, dw
\]
\textbf{Step 3: Spherical Polar Coordinates in $(u,v,w)$}
Let $u = \rho\sin\phi\cos\theta, v = \rho\sin\phi\sin\theta, w = \rho\cos\phi$.
For the positive octant: $\rho \in [0, 1], \phi \in [0, \pi/2], \theta \in [0, \pi/2]$:
\[
\iiint_{V'} u v w \, du \, dv \, dw = \int_0^{\pi/2} \cos\theta\sin\theta \, d\theta \int_0^{\pi/2} \sin^3\phi\cos\phi \, d\phi \int_0^1 \rho^5 \, d\rho
\]
\begin{itemize}[leftmargin=6mm]
    \item $\int_0^{\pi/2} \sin\theta\cos\theta \, d\theta = \left[ \frac{\sin^2\theta}{2} \right]_0^{\pi/2} = \frac{1}{2}$
    \item $\int_0^{\pi/2} \sin^3\phi\cos\phi \, d\phi = \left[ \frac{\sin^4\phi}{4} \right]_0^{\pi/2} = \frac{1}{4}$
    \item $\int_0^1 \rho^5 \, d\rho = \frac{1}{6}$
\end{itemize}
Multiplying:
\[
\frac{1}{2} \times \frac{1}{4} \times \frac{1}{6} = \frac{1}{48}
\]
Therefore:
\[
\mathbf{I = \frac{a^2 b^2 c^2}{48}}
\]

\begin{markingrubric}
\textbf{Mark Distribution:}
$\bullet$ Ellipsoid transformation and Jacobian $abc$: \textbf{3 Marks}\\
$\bullet$ Spherical coordinate conversion for positive octant: \textbf{3 Marks}\\
$\bullet$ Evaluation of individual angular and radial integrals: \textbf{3 Marks}\\
$\bullet$ Final answer $a^2 b^2 c^2 / 48$: \textbf{1 Mark}
\end{markingrubric}
\end{solutionbox}

% Solution 36
\begin{solutionbox}[Solution to Question 36 (10 Marks)]
\textbf{Step 1: Statement of Green's Theorem}
\[
\oint_C (P dx + Q dy) = \iint_R \left(\frac{\partial Q}{\partial x} - \frac{\partial P}{\partial y}\right) dx \, dy
\]
Here $P = xy + y^2$ and $Q = x^2$.
Boundaries: $y = x^2$ and $y = x$.
Intersections: $x^2 = x \implies x = 0$ and $x = 1$. Vertices: $(0,0)$ and $(1,1)$.

\textbf{Step 2: Double Integral Evaluation}
\[
\frac{\partial Q}{\partial x} = 2x, \quad \frac{\partial P}{\partial y} = x + 2y
\]
\[
\frac{\partial Q}{\partial x} - \frac{\partial P}{\partial y} = 2x - (x + 2y) = x - 2y
\]
Set up the double integral:
\[
\iint_R (x - 2y) \, dy \, dx = \int_0^1 \int_{x^2}^x (x - 2y) \, dy \, dx = \int_0^1 \left[ xy - y^2 \right]_{x^2}^x dx
\]
\[
= \int_0^1 \left[ (x^2 - x^2) - (x^3 - x^4) \right] dx = \int_0^1 (x^4 - x^3) dx = \left[ \frac{x^5}{5} - \frac{x^4}{4} \right]_0^1 = \frac{1}{5} - \frac{1}{4} = \mathbf{-\frac{1}{20}}
\]

\textbf{Step 3: Line Integral Evaluation ($\oint_C P dx + Q dy$)}
The curve $C = C_1 \cup C_2$:
\begin{enumerate}[label=(\roman*)]
    \item Along $C_1: y = x^2$ from $(0,0)$ to $(1,1)$, $dy = 2x dx$:
    \[
    \int_{C_1} = \int_0^1 \left[ (x(x^2) + (x^2)^2) dx + x^2(2x dx) \right] = \int_0^1 (x^3 + x^4 + 2x^3) dx = \int_0^1 (3x^3 + x^4) dx = \frac{3}{4} + \frac{1}{5} = \frac{19}{20}
    \]
    \item Along $C_2: y = x$ from $(1,1)$ to $(0,0)$, $dy = dx$:
    \[
    \int_{C_2} = \int_1^0 \left[ (x^2 + x^2) dx + x^2 dx \right] = \int_1^0 3x^2 dx = [x^3]_1^0 = 0 - 1 = -1
    \]
\end{enumerate}
Total line integral:
\[
\oint_C = \int_{C_1} + \int_{C_2} = \frac{19}{20} - 1 = \mathbf{-\frac{1}{20}}
\]
Since Line Integral ($-1/20$) = Double Integral ($-1/20$), **Green's Theorem is verified!**

\begin{markingrubric}
\textbf{Mark Distribution:}
$\bullet$ Stating Green's Theorem and computing $\frac{\partial Q}{\partial x} - \frac{\partial P}{\partial y}$: \textbf{2 Marks}\\
$\bullet$ Evaluating double integral yielding $-1/20$: \textbf{3 Marks}\\
$\bullet$ Evaluating line integrals along $C_1$ and $C_2$: \textbf{4 Marks}\\
$\bullet$ Equating results and concluding verification: \textbf{1 Mark}
\end{markingrubric}
\end{solutionbox}

\newpage
